\section{Noetherian schemes}
\label{section-noetherian}

\noindent
Recall that a ring $R$ is {\it Noetherian} if it satisfies the ascending
chain condition of ideals. Equivalently every ideal of $R$ is finitely
generated.

\begin{definition}
\label{definition-noetherian}
Let $X$ be a scheme.
\begin{enumerate}
\item We say $X$ is {\it locally Noetherian} if every
$x \in X$ has an affine open neighbourhood
$\Spec(R) = U \subset X$ such that the ring $R$ is Noetherian.
\item We say $X$ is {\it Noetherian} if $X$ is locally Noetherian
and quasi-compact.
\end{enumerate}
\end{definition}

\noindent
Here is the standard result characterizing locally Noetherian schemes.

\begin{lemma}
\label{lemma-locally-Noetherian}
Let $X$ be a scheme. The following are equivalent:
\begin{enumerate}
\item The scheme $X$ is locally Noetherian.
\item For every affine open $U \subset X$ the ring $\mathcal{O}_X(U)$
is Noetherian.
\item There exists an affine open covering $X = \bigcup U_i$ such that
each $\mathcal{O}_X(U_i)$ is Noetherian.
\item There exists an open covering $X = \bigcup X_j$
such that each open subscheme $X_j$ is locally Noetherian.
\end{enumerate}
Moreover, if $X$ is locally Noetherian then every open subscheme
is locally Noetherian.
\end{lemma}

\begin{proof}
To show this it suffices to show that being Noetherian is a local
property of rings, see Lemma \ref{lemma-locally-P}.
Any localization of a Noetherian ring is Noetherian, see
Algebra, Lemma \ref{algebra-lemma-Noetherian-permanence}.
By Algebra, Lemma \ref{algebra-lemma-cover} we see the second
property to Definition \ref{definition-property-local}.
\end{proof}

\begin{lemma}
\label{lemma-immersion-into-noetherian}
Any immersion $Z \to X$ with $X$ locally Noetherian is quasi-compact.
\end{lemma}

\begin{proof}
A closed immersion is clearly quasi-compact.
A composition of quasi-compact morphisms is quasi-compact,
see Topology, Lemma \ref{topology-lemma-composition-quasi-compact}.
Hence it suffices to show that an open immersion into
a locally Noetherian scheme is quasi-compact.
Using Schemes, Lemma \ref{schemes-lemma-quasi-compact-affine}
we reduce to the case where $X$ is affine.
Any open subset of the spectrum of a Noetherian ring
is quasi-compact (for example
combine Algebra, Lemma \ref{algebra-lemma-Noetherian-topology} and
Topology, Lemmas \ref{topology-lemma-Noetherian} and
\ref{topology-lemma-Noetherian-quasi-compact}).
\end{proof}

\begin{lemma}
\label{lemma-locally-Noetherian-quasi-separated}
A locally Noetherian scheme is quasi-separated.
\end{lemma}

\begin{proof}
By Schemes, Lemma \ref{schemes-lemma-characterize-quasi-separated}
we have to show that the intersection $U \cap V$ of two
affine opens of $X$ is quasi-compact. This follows from
Lemma \ref{lemma-immersion-into-noetherian} above on
considering the open immersion $U \cap V \to U$ for example.
(But really it is just because any open of the spectrum of a
Noetherian ring is quasi-compact.)
\end{proof}

\begin{lemma}
\label{lemma-Noetherian-topology}
A (locally) Noetherian scheme has a (locally)
Noetherian underlying topological space,
see Topology, Definition \ref{topology-definition-noetherian}.
\end{lemma}

\begin{proof}
This is because a Noetherian scheme is a finite union of spectra
of Noetherian rings and
Algebra, Lemma \ref{algebra-lemma-Noetherian-topology} and
Topology, Lemma \ref{topology-lemma-finite-union-Noetherian}.
\end{proof}

\begin{lemma}
\label{lemma-locally-closed-in-Noetherian}
Any locally closed subscheme of a (locally) Noetherian
scheme is (locally) Noetherian.
\end{lemma}

\begin{proof}
Omitted. Hint: Any quotient, and any localization of a Noetherian
ring is Noetherian. For the Noetherian case use again
that any subset of a Noetherian space is a Noetherian space
(with induced topology).
\end{proof}

\begin{lemma}
\label{lemma-Noetherian-irreducible-components}
A Noetherian scheme has a finite number of irreducible components.
\end{lemma}

\begin{proof}
The underlying topological space of a Noetherian scheme is Noetherian
(Lemma \ref{lemma-Noetherian-topology})
and we conclude because a Noetherian topological space
has only finitely many irreducible components
(Topology, Lemma \ref{topology-lemma-Noetherian}).
\end{proof}

\begin{lemma}
\label{lemma-morphism-Noetherian-schemes-quasi-compact}
Any morphism of schemes $f : X \to Y$ with $X$ Noetherian
is quasi-compact.
\end{lemma}

\begin{proof}
Use Lemma \ref{lemma-Noetherian-topology}
and use that any subset of a Noetherian topological
space is quasi-compact (see Topology,
Lemmas \ref{topology-lemma-Noetherian} and
\ref{topology-lemma-Noetherian-quasi-compact}).
\end{proof}

\noindent
Here is a fun lemma.
It says that every locally Noetherian scheme has plenty of
closed points (at least one in every closed subset).

\begin{lemma}
\label{lemma-locally-Noetherian-closed-point}
Any nonempty locally Noetherian scheme has a closed point.
Any nonempty closed subset of a locally Noetherian scheme has a closed point.
Equivalently, any point of a locally Noetherian scheme specializes
to a closed point.
\end{lemma}

\begin{proof}
The second assertion follows from the first (using
Schemes, Lemma \ref{schemes-lemma-reduced-closed-subscheme}
and Lemma \ref{lemma-locally-closed-in-Noetherian}).
Consider any nonempty affine open $U \subset X$.
Let $x \in U$ be a closed point. If $x$ is a closed point
of $X$ then we are done. If not, let $X_0 \subset X$ be the
reduced induced closed subscheme structure on $\overline{\{x\}}$.
Then $U_0 = U \cap X_0$ is an affine open of $X_0$ by
Schemes, Lemma \ref{schemes-lemma-closed-subspace-scheme} and
$U_0 = \{x\}$. Let $y \in X_0$, $y \not = x$ be a specialization of $x$.
Consider the local ring $R = \mathcal{O}_{X_0, y}$.
This is a Noetherian local ring as $X_0$ is Noetherian
by Lemma \ref{lemma-locally-closed-in-Noetherian}. Denote $V \subset \Spec(R)$
the inverse image of $U_0$ in $\Spec(R)$ by the canonical morphism
$\Spec(R) \to X_0$ (see Schemes, Section \ref{schemes-section-points}.)
By construction $V$ is a singleton with unique point corresponding to $x$ (use
Schemes, Lemma \ref{schemes-lemma-specialize-points}).
By
Algebra, Lemma \ref{algebra-lemma-Noetherian-local-domain-dim-2-infinite-opens}
we see that $\dim(R) = 1$.
In other words, we see that $y$ is an immediate specialization
of $x$ (see Topology, Definition \ref{topology-definition-dimension-function}).
In other words, any
point $y \not = x$ such that $x \leadsto y$ is an immediate
specialization of $x$. Clearly each of these points is a
closed point as desired.
\end{proof}

\begin{lemma}
\label{lemma-locally-Noetherian-specialization-dvr}
Let $X$ be a locally Noetherian scheme.
Let $x' \leadsto x$ be a specialization of points of $X$. Then
\begin{enumerate}
\item there exists a discrete valuation ring $R$ and a morphism
$f : \Spec(R) \to X$ such that the generic point $\eta$ of
$\Spec(R)$ maps to $x'$ and the special point maps to $x$, and
\item provided $x \not = x'$ given a finitely generated field
extension $K/\kappa(x')$, we may arrange it so that the extension
$\kappa(\eta)/\kappa(x')$
induced by $f$ is isomorphic to the given one.
\end{enumerate}
\end{lemma}

\begin{proof}
First assume $x' \leadsto x$ is a specialization in $X$
with $x \not = x'$ and let $K/\kappa(x')$ be a finitely
generated extension of fields. By
Schemes, Lemma \ref{schemes-lemma-specialize-points}
and the discussion following
Schemes, Lemma \ref{schemes-lemma-characterize-points}
this leads to ring maps $\mathcal{O}_{X, x} \to \kappa(x') \to K$.
Since $x \not = x'$ the image of $\mathcal{O}_{X, x}$ in $K$
is not a field (details omitted).
Let $R \subset K$ be any discrete valuation ring whose field of fractions is
$K$ and which dominates the image of $\mathcal{O}_{X, x} \to K$, see
Algebra, Lemma \ref{algebra-lemma-exists-dvr}.
The ring map $\mathcal{O}_{X, x} \to R$ induces the morphism
$f : \Spec(R) \to X$, see
Schemes, Lemma \ref{schemes-lemma-morphism-from-spec-local-ring}.
This morphism has all the desired properties by construction.
If $x = x'$ then we can set $R = \kappa(x)[t]_{(t)}$.
\end{proof}

\begin{lemma}
\label{lemma-thin-infinite-sequence}
Let $S$ be a Noetherian scheme. Let $T \subset S$ be an infinite subset.
Then there exists an infinite subset $T' \subset T$
such that there are no nontrivial specializations among the points $T'$.
\end{lemma}

\begin{proof}
Let $T_0 \subset T$ be the set of $t \in T$ which do not specialize
to another point of $T$. If $T_0$ is infinite, then $T' = T_0$ works.
Hence we may and do assume $T_0$ is finite.
Inductively, for $i > 0$, consider the set $T_i \subset T$
of $t \in T$ such that
\begin{enumerate}
\item $t \not \in T_{i - 1} \cup T_{i - 2} \cup \ldots \cup T_0$,
\item there exist a nontrivial specialization $t \leadsto t'$ with
$t' \in T_{i - 1}$, and
\item for any nontrivial specialization
$t \leadsto t'$ with $t' \in T$ we have
$t' \in T_{i - 1} \cup T_{i - 2} \cup \ldots \cup T_0$.
\end{enumerate}
Again, if $T_i$ is infinite, then $T' = T_i$ works.
Let $d$ be the maximum of the dimensions of the local rings
$\mathcal{O}_{S, t}$ for $t \in T_0$; then $d$ is an integer
because $T_0$ is finite and the dimensions of the local rings
are finite by Algebra, Proposition \ref{algebra-proposition-dimension}.
Then $T_i = \emptyset$ for $i > d$.
Namely, if $t \in T_i$ then we can find a sequence
of nontrivial specializations
$t = t_i \leadsto t_{i - 1} \leadsto \ldots \leadsto t_0$
with $t_0 \in T_0$. As
the points $t = t_i, t_{i - 1}, \ldots, t_0$ are in
$\Spec(\mathcal{O}_{S, t_0})$
(Schemes, Lemma \ref{schemes-lemma-specialize-points}),
we see that $i \leq d$.
Thus $\bigcup T_i = T_d \cup \ldots \cup T_0$ is a finite subset of $T$.

\medskip\noindent
Suppose $t \in T$ is not in $\bigcup T_i$. Then there must be a specialization
$t \leadsto t'$ with $t' \in T$ and $t' \not \in \bigcup T_i$. (Namely, if
every specialization of $t$ is in the finite set $T_d \cup \ldots \cup T_0$,
then there is a maximum $i$ such that there is some specialization
$t \leadsto t'$ with $t' \in T_i$ and then $t \in T_{i + 1}$ by construction.)
Hence we get an infinite sequence
$$
t \leadsto t' \leadsto t'' \leadsto \ldots
$$
of nontrivial specializations between points of $T \setminus \bigcup T_i$.
This is impossible because the underlying topological space of $S$
is Noetherian by Lemma \ref{lemma-locally-Noetherian-quasi-separated}.
\end{proof}

\begin{lemma}
\label{lemma-maximal-points}
Let $S$ be a Noetherian scheme. Let $T \subset S$ be a subset. Let
$T_0 \subset T$ be the set of $t \in T$ such that there is no nontrivial
specialization $t' \leadsto t$ with $t' \in T'$. Then (a) there are
no specializations among the points of $T_0$, (b) every point of
$T$ is a specialization of a point of $T_0$, and (c) the closures
of $T$ and $T_0$ are the same.
\end{lemma}

\begin{proof}
Recall that $\dim(\mathcal{O}_{S, s}) < \infty$ for any $s \in S$, see
Algebra, Proposition \ref{algebra-proposition-dimension}. Let $t \in T$.
If $t' \leadsto t$, then by dimension theory
$\dim(\mathcal{O}_{S, t'}) \leq \dim(\mathcal{O}_{S, t})$
with equality if and only if $t' = t$. Thus if we pick $t' \leadsto t$
with $\dim(\mathcal{O}_{T, t'})$ minimal, then $t' \in T_0$.
In other words, 
every $t \in T$ is the specialization of an element of $T_0$.
\end{proof}

\begin{lemma}
\label{lemma-countable-dense-subset}
Let $S$ be a Noetherian scheme. Let $T \subset S$ be an infinite dense subset.
Then there exist a countable subset $E \subset T$ which is dense in $S$.
\end{lemma}

\begin{proof}
Let $T'$ be the set of points $s \in S$ such that $\overline{\{s\}} \cap T$
contains a countable subset whose closure is $\overline{\{s\}}$.
Since a finite set is countable we have $T \subset T'$.
For $s \in T'$ choose such a countable subset
$E_s \subset \overline{\{s\}} \cap T$.
Let $E' = \{s_1, s_2, s_3, \ldots\} \subset T'$
be a countable subset. Then the closure of $E'$ in $S$ is the
closure of the countable subset $\bigcup_n E_{s_n}$ of $T$.
It follows that if $Z$
is an irreducible component of the closure of $E'$, then the generic
point of $Z$ is in $T'$.

\medskip\noindent
Denote $T'_0 \subset T'$ the subset of $t \in T'$ such that
there is no nontrivial specialization $t' \leadsto t$ with $t' \in T'$
as in Lemma \ref{lemma-maximal-points} whose results we will use
without further mention. If $T'_0$ is infinite, then we choose a
countable subset $E' \subset T'_0$. By the argument in the first
paragraph, the generic points of the irreducible components of the
closure of $E'$ are in $T'$. However, since one of these points specializes to
infinitely many distinct elements of $E' \subset T'_0$
this is a contradiction. Thus $T'_0$ is finite, say
$T'_0 = \{s_1, \ldots, s_m\}$. Then it follows that $S$, which is
the closure of $T$, is contained in the closure of
$\{s_1, \ldots, s_m\}$, which in turn is contained in the closure
of the countable subset $E_{s_1} \cup \ldots \cup E_{s_m} \subset T$
as desired.
\end{proof}








